% ---
% id: protocol-compiler/rem:gate-by-certificate-form
% kind: rem
% title: "The source gate on each certificate form"
% summary: >-
%   Gate (G2) refuses every block-form certificate by its form, pure or not,
%   and admits an ensemble presentation, whose blocks are rank one by theorem;
%   the refusal rests on scope, not on that difference. At several Alice
%   sectors a presentation is not globally pure, and the gate does not demand it.
% prominence: standard
% status: accepted
% uses: [def:eligible]
% mentions: [def:declaration, thm:gates, transmitter/rem:certificate-forms, transmitter/thm:certificate-sufficiency, transmitter/lem:ensemble-presentation]
% assumes: []
% refs: []
% provenance: newly-derived
% lean:
% ---

Gate (G2) of Definition~\ref{def:eligible} re-decides a pure-form
certificate, refuses a block-form one and decides an ensemble presentation by
what it declares. What purity adds to a certificate, and why the certificate
has a block form and an ensemble presentation at all, is
\xref{transmitter/rem:certificate-forms}; this remark says what separates the
forms at the gate.

\begin{remark}\label{rem:gate-by-certificate-form}
\emph{The block form.} Gate (G2) decides $\Psi^2=\Psi$ of a pure-form
certificate and refuses a \tref[block form]{block-form} certificate by its
declared form, whatever its blocks (Definition~\ref{def:eligible}). Most
block-form certificates fail $\Psi^2=\Psi$ as well --- every one with two or
more non-zero blocks, and every one whose one non-zero block is not a
rank-one projector --- but one whose one non-zero block is a rank-one
projector is pure and is refused all the same. So (G2) is not the consumer that
\xref{transmitter/rem:certificate-forms} describes, which demands purity of
every certificate and refuses a block form for want of it: the ground of the
refusal is the scope rule stated at the end of this remark.

\emph{What separates a general block form from an ensemble presentation} is
not the word ``purity'', which is waived on both sides once Alice carries
several sectors. A general block form declares blocks that are positive
semidefinite and \emph{not known to be rank one}, so the ``only if'' half of
\xref{transmitter/thm:certificate-sufficiency}\,(ii) --- proved for a block
form exactly where every non-zero block is a positive multiple of a rank-one
projector --- is lost with it, and a source/POVM mismatch would then be
excluded only by the POVM condition decided on its own. A presentation's
blocks are rank one \emph{by theorem}, so that half is not lost.

\emph{An ensemble presentation at one Alice sector.} For an \tref{ensemble
presentation} there is no declared $\Psi$ and no declared $\{M_x\}$ to
re-decide: purity of the blocks, (SR1) and (SR2) are \emph{theorems} of the
presentation (\xref{transmitter/lem:ensemble-presentation}\,(i), (ii)), that
lemma's (ii) giving every non-zero block as a positive multiple of a rank-one
projector. At one Alice sector no purity requirement is waived: there the
$\Psi$ the presentation defines is that one block, of unit trace, hence a
rank-one projector, so the five conditions of (G2) hold of it as theorems
rather than as decisions, and the gate demands of such a presentation
everything it demands of a pure-form certificate.

\emph{At two or more Alice sectors.} There the defined operator is
$\Psi=\bigoplus_t\Psi_t$ with $\tr\Psi_t=g_t$ and $\sum_tg_t=1$, so
$\Psi^2=\bigoplus_tg_t\Psi_t\ne\Psi$ as soon as two blocks are non-zero, by
the same arithmetic by which a block form with two non-zero blocks fails
$\Psi^2=\Psi$. Global purity is therefore not
demanded of a presentation with several Alice sectors. What the gate rests on
there instead is two statements of the transmitter's, both available for a
presentation, of which a general block form has only the second.
\begin{enumerate}[label=(\alph*),leftmargin=2.2em]
\item Every block is rank one by theorem
  (\xref{transmitter/lem:ensemble-presentation}\,(ii)), which is exactly the
  hypothesis under which \xref{transmitter/thm:certificate-sufficiency}\,(ii)
  holds in \emph{both} directions.
\item \xref{transmitter/thm:certificate-sufficiency}\,(iii), the sufficiency
  statement at the \tref{module boundary}, assumes (SR1)--(SR3) and that
  module's sequential assumption and \emph{reads purity nowhere}.
\end{enumerate}
The asymmetry in (a) is not relied on. The ``only if'' direction is what a
consumer that does \emph{not} decide the POVM condition separately would need,
and this module decides it at check V3 in either form ---
$\sum_xM^{(t)}_x=\idon{A_t}$ in \tref{POVM form}, the ensemble identity of
\eqref{eq:ensembleform} in ensemble form (D(iii)) --- which is stronger than
the compressed condition that converse delivers; and for a presentation there
is in any case no mismatch to detect, $\{M_x\}$ being \emph{defined} by the
presentation rather than declared and gate (G2) pinning the declaration's
Alice slots to the transposed declared blocks entry for entry, so that
direction holds vacuously. Clause (a) records what the two forms differ in;
it does not carry the refusal of a general block form, which rests on scope:
admitting one would widen (G2) in the unsafe direction, to certificates that
can accompany a declaration in the \tref{basic format}.

\emph{Where a presentation can occur.} A presentation is admitted only beside
an ensemble-form declaration, and no declaration in the basic format is in
ensemble form.

\emph{No form's (SR1)--(SR2) is assumed.} In none of its forms are the
conditions (SR1)--(SR2) of a \tref{source-replacement certificate} an
assumption of this module: what gate (G2) decides of each form is
Definition~\ref{def:eligible}'s and Theorem~\ref{thm:gates}'s, and
(SR1)--(SR2) are named as a hypothesis at the one place a derivation uses
them, \xref{transmitter/thm:certificate-sufficiency}.
\end{remark}
